\section{Related Work}
\label{related}

\noindent\textit{Synchronous VSS-Based Discrete-Logarithm Constructions.}
Pedersen's DKG~\cite{Pedersen91} is a two-round protocol among $n$ parties,
such that each party plays the role of the dealer for Feldman Verifiable Secret
Sharing (VSS)~\cite{Feldman87}.
Gennaro et al.~\cite{GennaroJKR07} describe an attack against Pedersen's DKG, which we refer to as the \GJKRAttack attack,
that considers a \emph{rushing adversary}, who is allowed to speak last.
In this attack,
an adversary can influence the distribution of key material
by adaptively choosing their contribution.
While the authors later demonstrate unforgeability of a Schnorr threshold signature scheme with  Pedersen's DKG employed for key generation~\cite{GennaroJK02},
this proof of security does \emph{not} extend to other applications.
%
Gurkan et al.~\cite{GurkanJMMST21} show that Pedersen's DKG can be used securely in any setting where the cryptographic scheme is ``rekeyable''.
In other words,
if the scheme produces some cryptographic output
$o_1$ (such as a ciphertext or signature) that is valid with respect to a secret key $\sk_1$,
it is possible to obtain a second value $o_2 \gets f_1(o_1, \sk', \alpha)$,
such that $o_2$ is valid with respect to a related secret key $\sk_2 \gets f_2(\sk_1, \sk', \alpha)$,
where $\alpha$ is a side input and $f_1, f_2$ are efficient update functions.
While some constructions such as BLS signatures~\cite{BonehLS04} do support rekeyability,
others such as Schnorr signatures do not.

Gennaro et al.~\cite{GennaroJKR07} present a construction secure against the
\GJKRAttack attack
that is (similar to our generic construction) three rounds in the optimistic setting,
but it has higher round complexity when parties misbehave.
While their description of their construction implies that a round to broadcast failure messages is optional,
their construction is in fact insecure without this round,
as a malicious party may force a subset of parties to fail,
and so selectively choose which participant contributions are used towards the
creation of the output public key and secret key shares
(thereby biasing the output key material).
While requiring higher round complexity,
their construction offers strong robustness as players negotiate in the second
and third round which qualified set of players to continue the protocol with.
Our generic construction simply aborts the protocol,
requiring the implementation to investigate the set of cheating players outside
the protocol, a reasonable assumption in practice.

Lindell~\cite{Lindell22} introduced a DKG that builds
upon Schnorr proofs of knowledge and online-extractable zero-knowledge proofs of knowledge.
However, the protocol incurs computational costs due to the
requirement of online extractors~\cite{Fischlin05}.
Furthermore, the security of the scheme is only generally discussed without reference to a concrete notion of security.


Many efficient pairing-based DKGs and secret-sharing primitives are described in the
literature~\cite{GurkanJMMST21,Groth21,AbrahamJMMS2022,ZhangXHSZ22}.
Further, concrete DKGs have been presented in the ECDSA and RSA
settings~\cite{LindellN18,DamgardM10,FrankelMY98,CanettiGGMP20}.
This work presents a generic construction that can be instantiated
without pairings.


\medskip

\noindent\textit{Synchronous PVSS-Based Discrete-Logarithm Constructions.}
Our generic constructions assume VSS as a building block,
and consequently incur an extra round of communication.
It is possible to instead use a publicly verifiable secret sharing scheme (PVSS),
as is demonstrated by alternative DKGs in the literature.
However,
PVSS schemes require that secret key material be encrypted to receivers' public keys,
thereby requiring a complicated decryption step.
For this reason,
we define our construction to explicitly incur an additional network round.
However, our construction could instead be adapted to use a PVSS and thereby
use only two network rounds.


\medskip

\noindent\textit{Definitions of DKG Security.}
Gurkan et al.~\cite{GurkanJMMST21} present a definition of DKG security, in which the DKG is secure if its use within some larger protocol
does not weaken the security of the overall scheme from the centralized (single-party) setting.
We build upon this notion by defining the security of a DKG with respect to a target key generation scheme.
%
Bacho and Loss~\cite{BachoL22}
introduce the security notions of consistency and oracle-aided algebraic
simulatability for a DKG.
Their notion of consistency is encompassed by our notions of strong and weak robustness.
Because their notion of oracle-aided simulatability assumes a discrete logarithm oracle, it is not applicable to schemes without such algebraic structure.
%%
Boneh and Shoup~\cite{BonehS2020} present a generic definition for the security
of a distributed key generation protocol using a simulation-based approach,
with respect to a target trusted secret sharing scheme.
However, their definition is given in a UC-based approach,
whereas we present security experiments using a game-based approach to
explicitly define subtle details such as participation communication patterns and aborts.
However, our notions do not guarantee security under composition with other schemes.
